\begin{theorem}[Norm filtration above the break]
\label{mod:ramification-normabovebreak}
Let $K/F$ be a cyclic extension of prime degree
$\ell=[K:F]$, let $t\in\mathbf N$ be its lower ramification break, and
assume $f(K/F)=1$.  Fix an integral uniformizer $\pi_K$ satisfying the
explicit monogenicity hypothesis
\[
 \operatorname{Algebra.adjoin}_{\mathcal O_F}\{\pi_K\}=\mathcal O_K.
\]
For a natural target depth $r$ with the strict inequality $t<r$, Lean uses
the manuscript's one-break Herbrand depth
\[
 a=\texttt{aboveBreakSourceDepth}(t,\ell,r)
   =\Psi_{t,\ell}(r)=t+\ell(r-t).
\]
In particular, the theorem is not stated at the critical depth $r=t$.

Put $D=(\ell-1)(t+1)$.  The two integer-division identities needed at the
adjacent source layers are proved separately:
\[
 \left\lfloor\frac{a+D}{\ell}\right\rfloor=r,
 \qquad
 \left\lfloor\frac{a+1+D}{\ell}\right\rfloor=r+1.
\]
After substituting $e(K/F)=\ell$, which follows from $f(K/F)=1$, and the
different formula of module~\ref{mod:ramification-different}, the exact
trace-image theorem of module~\ref{mod:ramification-traceideals} therefore
gives
\[
 \operatorname{Tr}_{K/F}(\mathfrak p_K^a)=\mathfrak p_F^r,
 \qquad
 \operatorname{Tr}_{K/F}(\mathfrak p_K^{a+1})=\mathfrak p_F^{r+1}.
\]
These are equalities of mapped submodules; the first is also used in its
elementwise surjectivity form when constructing graded norm preimages.

For every $x\in K$, Lean records the exact expansion
\[
 N_{K/F}(1+x)-1=\sum_{j=1}^{\ell}s_j(x),
\]
where $s_j$ is the $j$th elementary symmetric function of the conjugates
of $x$.  If $v_K(x)\ge a$, then every higher term, including the terminal
norm coefficient, satisfies the exact next-layer bound
\[
 s_j(x)\in\mathfrak p_F^{r+1}
 \qquad(2\le j\le\ell).
\]
For $2\le j<\ell$, the positive-break branch derives
$\operatorname{char}k_F=\ell$ from the faithful ramification coordinate
and applies the strong wild symmetric bound; the $t=0$ branch applies the
totally ramified ceiling bound.  The terminal case $j=\ell$ uses
$e(K/F)=\ell$ and the exact valuation of the norm.  Thus no coefficient is
lost at an endpoint, and the expansion gives the precise congruence
\[
 N_{K/F}(1+x)-1-\operatorname{Tr}_{K/F}(x)
   \in\mathfrak p_F^{r+1}.
\]

The predicate \texttt{AboveBreakNormMapsUnitFiltration} records a certified
containment between a source and target unit layer.  Lean proves the two
distinct containments
\[
 N_{K/F}(U_K^a)\subseteq U_F^r,
 \qquad
 N_{K/F}(U_K^{a+1})\subseteq U_F^{r+1},
\]
and only then descends norm to the heterogeneous quotient map
\[
 \overline N_r:U_K^a/U_K^{a+1}\longrightarrow U_F^r/U_F^{r+1}.
\]
The numerator and denominator containments are separate arguments, so the
map is independent of representatives by construction.  On a certified
representative $u=1+x$ its quotient-class formula is exactly
\[
 [1+x]\longmapsto[1+\operatorname{Tr}_{K/F}(x)].
\]
The general quotient construction also proves that norm commutes with the
canonical projections obtained by deepening both denominator layers, as
well as the representative formula and the fact that a representative in
the deeper source layer maps to the identity class.

Exact trace surjectivity and the displayed congruence show that
$\overline N_r$ is surjective at every $r>t$.  Since $f(K/F)=1$ identifies
the finite residue cardinalities, the source and target positive graded
pieces have the same order; hence the graded norm is bijective and Lean
packages it as the multiplicative equivalence
\[
 U_K^{\Psi_{t,\ell}(r)}/U_K^{\Psi_{t,\ell}(r)+1}
   \simeq U_F^r/U_F^{r+1}.
\]

Finally, Lean applies complete successive lifting to the continuous norm.
For $a=\Psi_{t,\ell}(r)$ it passes the predecessor depths
\[
 b(n)=(a-1)+\ell n,
 \qquad b(n)+1=\Psi_{t,\ell}(r+n),
\]
to module~\ref{mod:localfield-successivelifting}.  These depths are
monotone and tend to infinity, while the local-field hypotheses supply
completeness and continuity.  The resulting full-layer surjectivity is
stated elementwise and then as the exact subgroup image
\[
 \boxed{N_{K/F}(U_K^{\Psi_{t,\ell}(r)})=U_F^r}.
\]
For the adjacent source layer Lean does not identify
$\Psi_{t,\ell}(r)+1$ with $\Psi_{t,\ell}(r+1)$: it uses
\[
 \Psi_{t,\ell}(r)+1<\Psi_{t,\ell}(r+1)
   =\Psi_{t,\ell}(r)+\ell,
\]
combines surjectivity from the deeper Herbrand layer with the separately
proved denominator containment, and obtains the second exact image
\[
 \boxed{N_{K/F}(U_K^{\Psi_{t,\ell}(r)+1})=U_F^{r+1}}.
\]
Both formulas are also exported as elementwise if-and-only-if statements.

\LeanFile{LanglandsFirstMainLemma/Ramification/NormAboveBreak.lean}
\lean{LanglandsFirstMainLemma.aboveBreakSourceDepth}
\lean{LanglandsFirstMainLemma.aboveBreakSourceDepth_eq}
\lean{LanglandsFirstMainLemma.aboveBreak_trace_floor}
\lean{LanglandsFirstMainLemma.aboveBreak_trace_floor_succ}
\lean{LanglandsFirstMainLemma.residueCharacteristic_eq_degree_of_positive_break}
\lean{LanglandsFirstMainLemma.norm_one_add_sub_one_eq_sum_elementarySymmetric}
\lean{LanglandsFirstMainLemma.higher_elementarySymmetric_mem_lattice_aboveBreakCanonical}
\lean{LanglandsFirstMainLemma.norm_one_add_sub_one_sub_trace_mem_aboveBreakCanonical}
\lean{LanglandsFirstMainLemma.trace_lattice_image_aboveBreak}
\lean{LanglandsFirstMainLemma.trace_lattice_image_aboveBreak_succ}
\lean{LanglandsFirstMainLemma.AboveBreakNormMapsUnitFiltration}
\lean{LanglandsFirstMainLemma.aboveBreakNormUnitFiltrationHom}
\lean{LanglandsFirstMainLemma.aboveBreakNormUnitFiltrationQuotient}
\lean{LanglandsFirstMainLemma.aboveBreakNormUnitFiltrationQuotient_mk}
\lean{LanglandsFirstMainLemma.aboveBreakNormUnitFiltrationQuotient_deeper_eq_one}
\lean{LanglandsFirstMainLemma.aboveBreakNormUnitFiltrationQuotient_projection}
\lean{LanglandsFirstMainLemma.normMapsUnitFiltration_aboveBreakCanonical}
\lean{LanglandsFirstMainLemma.normMapsUnitFiltration_aboveBreakCanonical_succ}
\lean{LanglandsFirstMainLemma.normGradedAboveBreakCanonical}
\lean{LanglandsFirstMainLemma.normGradedAboveBreakCanonical_mk}
\lean{LanglandsFirstMainLemma.normGradedAboveBreakCanonical_mk_eq_trace}
\lean{LanglandsFirstMainLemma.norm_graded_surjective_above_break}
\lean{LanglandsFirstMainLemma.norm_graded_bijective_above_break}
\lean{LanglandsFirstMainLemma.normGradedAboveBreakEquiv}
\lean{LanglandsFirstMainLemma.normGradedAboveBreakEquiv_apply}
\lean{LanglandsFirstMainLemma.normGradedAboveBreakCanonical_range_eq_top}
\lean{LanglandsFirstMainLemma.norm_unitFiltration_surjective_aboveBreak}
\lean{LanglandsFirstMainLemma.mem_unitFiltration_iff_exists_norm_aboveBreak}
\lean{LanglandsFirstMainLemma.norm_unitFiltration_map_eq_aboveBreak}
\lean{LanglandsFirstMainLemma.mem_unitFiltration_succ_iff_exists_norm_aboveBreak}
\lean{LanglandsFirstMainLemma.norm_unitFiltration_map_eq_aboveBreak_succ}
\leanok
\SourceText{Theorem~\ref{thm:graded-norm}, parts (iii)--(iv),
Theorem~\ref{thm:norm-filtration}, part (d), and
Corollary~\ref{cor:wild-norm-truncation}.}
\uses{mod:ramification-herbrand,mod:ramification-traceideals,
mod:ramification-symmetricbounds,mod:localfield-successivelifting,
mod:localfield-finitequotients}
\FormalizationContract The strict hypothesis $t<r$, both $a$ and $a+1$
trace targets, every coefficient with $j\ge2$, the factors
$e(K/F)=\ell$ and $f(K/F)=1$, separate quotient well-definedness,
projection compatibility, graded bijectivity, complete lifting, and both
exact norm-image equalities are explicit.  The proof imports neither the
assembled norm-filtration theorem nor pullback conductors, and it never
replaces an equality by an inclusion.
\end{theorem}
